\documentclass[10pt,a4paper]{amsart}
\usepackage[margin=19mm]{geometry}
\usepackage{amsmath,amssymb,mathtools}
\usepackage[scr=rsfs,cal=euler]{mathalfa}
\usepackage{xfrac}
\usepackage[unicode,hidelinks,bookmarks=false]{hyperref}
\newcommand{\ie}{i.e.\ }
\newcommand{\cf}{cf.\ }
\newcommand{\resp}{respectively\ }
\newcommand{\Subsection}{\subsection}
\providecommand{\nobreakdash}{\nobreak\hbox{-}}
\setlength{\emergencystretch}{2em}
\multlinegap0pt
\usepackage{bm}
\usepackage{bbm}
\usepackage{dsfont}
\usepackage{xspace}
\usepackage{cancel}
\usepackage{mathtools}
\usepackage{amsthm}
\makeatletter
\@ifundefined{theorem}{\newtheorem{theorem}{Theorem}}{}
\makeatother
\title[Negative Moment Bounds for Sample Autocovariance Matrices of]{Negative Moment Bounds for Sample Autocovariance Matrices of Stationary Processes Driven by Conditional Heteroscedastic Errors and Their Applications}
\author{Hsueh-Han Huang, Ching-Kang Ing,, Shu-Hui Yu}
\date{Source: arXiv:2301.07476v6 (CC BY 4.0)}
\begin{document}
\maketitle

\noindent\emph{Source and licence.} Negative Moment Bounds for Sample Autocovariance Matrices of Stationary Processes Driven by Conditional Heteroscedastic Errors and Their Applications. Version arXiv:2301.07476v6 (CC BY 4.0), distributed under the Creative Commons Attribution 4.0 International licence, as recorded in the arXiv licence metadata for this exact version and shown on the abstract page https://arxiv.org/abs/2301.07476v6. Full rights notice: licenses/arxiv-2301.07476-CC-BY-4.0.txt.\par\smallskip

\noindent\textit{Editorial scope.} Counted unit: the Theorem whose source label is \texttt{tfisherS} in the article (source0.tex lines 2881--2911, source label \texttt{tfisherS}), delivered together with the article's own complete original proof of it (source0.tex lines 2913--2982, one \texttt{proof} environment of 70 lines). No mathematics is written, repaired or re-derived: statement and proof are the article's own text line for line. Required context retained uncounted: AR mean part (lines 129--131); stationary char poly (lines 160--162); t.fisher.inverse moment (lines 208--211); tfisher (lines 611--651); pt.fisher.3 (lines 2245--2253); p j less bar p j (lines 2455--2457); ing113 (lines 2548--2550). Every definition, display and label of those lines is retained unchanged. Omitted from this package and not redistributed: 1--128, 132--159, 163--207, 212--610, 652--2244, 2254--2454, 2458--2547, 2551--2880, 2912--2912, 2983--4320. Nothing inside a retained excerpt is omitted or altered: every retained body is delivered exactly as the article's lines are written, so no omission receipt of any kind accompanies this package. Citations: the retained excerpts carry no citation command at all, so no bibliography entry of the article is transcribed and no reference is redistributed. Reference adaptation: none was needed and none was made; every reference inside the delivered unit resolves inside it, so no unresolved reference or citation is produced. Printed slips kept literal: no printed word, symbol, hypothesis or number of the counted statement or of its proof was altered. Layout and class: the article's own class is \texttt{imsart} with packages including amsthm, amsmath, natbib, hyperref, graphicx, bm, algorithm, float, algpseudocode, accents, inputenc; the wrapper is the lane's pinned \texttt{amsart} editorial wrapper at 10pt on A4, which restates the theorem-like environments the retained text needs and adds the article's own macro definitions that the retained text needs (none), with extra wrapper packages (bm, bbm, dsfont, xspace, cancel, mathtools). The delivered numbering is the unit's own. Assets: no retained span contains a figure environment, a graphics-inclusion command, a PDF-page inclusion, a table or a TikZ picture; no image file is copied into this package, the package declares no asset dependency and no image file is redistributed with it. Licence: version arXiv:2301.07476v6 of this article is distributed under the Creative Commons Attribution 4.0 International licence, as recorded in the arXiv licence metadata for this exact version and shown on the abstract page https://arxiv.org/abs/2301.07476v6. A full rights notice is delivered at licenses/arxiv-2301.07476-CC-BY-4.0.txt. Characters: every retained excerpt is pure ASCII, so the delivered engine, XeLaTeX with its default Latin Modern text font, reports no missing character.
\begin{equation}\label{AR mean part}
x_t=\sum_{i=0}^{\infty}\alpha_i\varepsilon_{t-i},
\end{equation}

\begin{equation}\label{stationary char poly}
\sum_{i=0}^{\infty}\alpha_iz^i\neq0\ \textrm{for complex}\ |z|\leq1.
\end{equation}

\begin{eqnarray}
\label{t.fisher.inverse moment}
E[\lambda_{\min}^{-q}(\hat{\mathbf{R}}_n(k))]=O(1), \,\, \mbox{as}\,\, n \to \infty,
\end{eqnarray}

\begin{theorem}\label{tfisher}
Assume \eqref{AR mean part}--\eqref{stationary char poly}
and {\rm (CH)(i)}. 
Suppose for any 
%$t \in \mathbb{Z}$and 
$0< x< y<\infty$, %that $\phi_t(\cdot)$ is bounded above and symmetric about zero and obeys
\begin{align}\label{zt density}
\begin{split}
\phi_t(y)\leq \phi_t(x)\,\,\mbox{and}\,\,\phi_t(-y)\leq \phi_t(-x).
\end{split}
\end{align}
In addition, 
for any $0<\delta \leq 1$, there exist a
finite positive constant $M_\delta$
and a pair ($\bar{\theta} ,\bar{C} $), independent of $\delta$, satisfying $0<\bar{\theta} \leq 1$ and $0<\bar{C}<\infty$,
such that
\begin{eqnarray}
\label{ing104}
\sup_{-\infty<t<\infty}\{\phi_t(\delta)+\phi_t(-\delta)\}\leq M_\delta,\,\,\sup_{-\infty<t<\infty}
\int_{-\delta}^{\delta}\phi_t(x)dx\leq\bar{C}\delta^{\bar{\theta}}.
\end{eqnarray}
Moreover, for any $t\in \mathbb{Z}$ and $x \neq 0$, there exists
a positive number $m_t(x)$ obeying $c_1|x| < m_t(x)< c_2|x|$, 
%$c_1|x| < m^{+}_t(x), m_t^{-}(x)< c_2|x|$,
where $0<c_1<c_2<\infty$ are constants that do not depend on $t$ or $x$, such that
\begin{align}
\begin{split}
\label{zt density 0}
\zeta_{t,x}(m_t(x))=\sup_{c>0} \zeta_{t,x}(c),
\end{split}
\end{align}
and 
\begin{equation}
\label{zt density 2}
\begin{split}
\zeta_{t,x}(c) \,\,\mbox{is non-increasing for} \,\,
c>m_t(x).
\end{split}
\end{equation}
Then, \eqref{t.fisher.inverse moment} follows.
\end{theorem}

\begin{equation}\label{pt.fisher.3}
\begin{split}
     & E\{\lambda_{\min}^{-q}(\sum_{i=1}^{dk}\bm{\Phi}_{ik}\bm{\Phi}_{ik}^\top)\} =   \int_{0}^{\infty}P\{\lambda_{\min}(\sum_{i=1}^{dk}\bm{\Phi}_{ik}\bm{\Phi}_{ik}^\top)<u^{-\frac{1}{q}}\}\mathrm{d}u \\
    \leq & \tilde{M}+\int_{\tilde{M}}^{\infty}P\{\inf\limits_{\lVert \bm{y}\rVert=1}\sum_{i=1}^{dk}(\bm{\Phi}_{ik}^\top\bm{y})^2<u^{-\frac{1}{q}},
   \max\limits_{1\leq j\leq dk^2}\sigma_j^2<u^{\tilde{\theta}},\max\limits_{1\leq i\leq dk}\lVert\bm{\Phi}_{ik}\rVert<\frac{u^{\frac{l}{q}}}{\sqrt{k}}\}\mathrm{d}u \\
     & +\int_{\tilde{M}}^{\infty}P(\max\limits_{1\leq j\leq dk^2}\sigma_j^2\geq u^{\tilde{\theta}})du+\int_{\tilde{M}}^{\infty}P({\max\limits_{1\leq i\leq dk}\lVert\bm{\Phi}_{ik}\rVert\geq \frac{u^{\frac{l}{q}}}{\sqrt{k}}})\mathrm{d}u \\
    := & \tilde{M}+{\rm (I)+(II)+(III)}.
\end{split}
\end{equation}

\begin{equation}\label{p j less bar p j}
p_j(\mathbf{s})\leq\bar{p}_j(s_j),
\end{equation}

\begin{equation}\label{ing113}
\int_{B(u)}\phi_{(i-1)k+k^{\star}}(z)\mathrm{d}z\leq \bar{C}(k^{1/2}6u^{-1/(2q)}/c_0)^{\bar{\theta}}.
\end{equation}

\begin{theorem}\label{tfisherS}
{Assume \eqref{AR mean part}--\eqref{stationary char poly} and {\rm (CH)(i)}}. Suppose for all $t \in \mathbb{Z}$, there exist positive constants $\bar{c}_1$, $\bar{c}_2$, and $\bar{c}_3$ such that
\begin{align}\label{zt densityS}
\begin{split}
& \phi_t(y)\leq \phi_t(x),\ \textrm{for all}\  \bar{c}_1\leq x< y<\infty,\\
& \phi_t(-y)\leq \phi_{t}(-x),\ \textrm{for all}\  \bar{c}_2\leq x< y<\infty,
\end{split}
\end{align}
and
\begin{eqnarray}
\label{ing104S}
\sup_{-\infty<t,x<\infty}\phi_t(x)\leq \bar{c}_3.
\end{eqnarray}
Moreover, for any $t\in \mathbb{Z}$ and $x\neq0$, there exist a positive number $m_t(x)$ and positive constants $\bar{c}_4$ and $\bar{c}_5$ obeying $\bar{c}_4|x| < m_t(x)< \bar{c}_5|x|$, such that
\begin{align}
\begin{split}
\label{zt density 0S}
& \zeta_{t,x}\{m_t(x)\}=\sup_{c>0} \zeta_{t,x}(c),
\end{split}
\end{align}
and 
\begin{equation}
\label{zt density 2S}
\begin{split}
\zeta_{t,x}(c) \,\,\mbox{is non-increasing for} \,\,
c>\bar{c}_5|x|.
\end{split}
\end{equation}
Then,
\eqref{t.fisher.inverse moment} follows.
\end{theorem}

\begin{proof}[Proof of Theorem \ref{tfisherS}]
{
We continue using the notation as in the proof of
Theorem \ref{tfisher}.
In particular, we
abbreviate 
$p^{(ik,k)}_j(\cdot)$ and $\tilde{\sigma}^{(ik,k)}_j(\cdot)$ as $p_j(\cdot)$ and $\tilde{\sigma}_j(\cdot)$, respectively,
for
$j=(i-1)k+1,\ldots,ik$.
To complete this proof, it suffices to replace two earlier bounds, \eqref{p j less bar p j}
and \eqref{ing113} from the proof of Theorem \ref{tfisher},
with the following two new bounds:
\begin{equation}\label{ptfisherS.1}
p_j(\mathbf{s})\leq\tilde{p}_j(s_j), \,\,\mbox{for}\,\,u> (\bar{c}_1\vee\bar{c}_2)^{4/\tilde{\theta}} \vee
(c_0/\bar{c}_5)^{4/(3\tilde{\theta})} \,\,\mbox{and}\,\,
 \mathbf{s}\in A_{ik,k}(u),
%\ \textrm{for all}\ j=(i-1)k+1,\ldots,ik\ \textrm{and}\ \mathbf{s}\in A_{ik,k}(u),
\end{equation}
and 
\begin{equation}\label{ptfisherS.2}
\int_{B(\mathbf{s}^{\prime},\mathbf{s}^{\prime\prime},u)}\phi_{(i-1)k+k^{\star}}(z)\mathrm{d}z\leq \bar{c}_3(k^{1/2}6u^{-1/(2q)}/c_0),\,\,\mbox{for}\,\, u>0,
\end{equation}
respectively.
Here, $\tilde{p}_j(s_j)$ is defined as
\begin{equation*}
\tilde{p}_j(s_j)= 
  \begin{cases}
    \frac{1}{c_0}\phi_j(\frac{s_j}{c_0}) & \mbox{for } |s_j|<\frac{c_0}{\bar{c}_5}, \\
    \frac{\bar{c}_3}{\bar{c}_4|s_j|} & \mbox{for } \frac{c_0}{\bar{c}_5}\leq|s_j|\leq u^{\frac{3\tilde{\theta}}{4}}, \\
    \frac{1}{c_0}\phi_j(\frac{s_j}{u^{\frac{\tilde{\theta}}{2}}}) & \mbox{for } u^{\frac{3\tilde{\theta}}{4}}<|s_j|.
  \end{cases}
\end{equation*}
Note that
$\tilde{\theta}$ is defined right before \eqref{pt.fisher.3},
and for a given $1\leq i \leq dk$,
\eqref{ptfisherS.1}
is required to hold
for all
$j=(i-1)k+1,\ldots,ik$.
%hold under \eqref{zt densityS}--\eqref{zt density 2S}.
%Then, the remainder of the proof is the same as the proof of Theorem \ref{tfisher}. 
%Note that except for $\tilde{p}_j(\cdot)$, $\bar{c}_3$, $\bar{c}_4$, and $\bar{c}_5$, all other notations in \eqref{ptfisherS.1} and \eqref{ptfisherS.2} are consistent with those used in the proof of Theorem \ref{tfisher}. 

 
 Since we have shown in the proof of 
 Theorem \ref{tfisher} that the length of $B(\mathbf{s}^{\prime},\mathbf{s}^{\prime\prime},u)$ is bounded by $k^{1/2}6u^{-1/(2q)}/c_0$,
\eqref{ptfisherS.2} follows directly from
\eqref{ing104S}.
To establish
\eqref{ptfisherS.1}, first observe that the inequality holds trivially 
when $s_j=0$.
For $s_j \neq 0$ and
$|s_j|<c_0/\bar{c}_5$, \eqref{zt density 2S} implies that $(1/c)\phi_j(s_j/c)$ is non-increasing in $c$ for all $c\geq c_0$. Since $\tilde{\sigma}_j(\mathbf{s})\geq c_0$, it follows that $p_j(\mathbf{s})\leq(1/c_0)\phi_j(s_j/c_0)$. For the case $c_0/\bar{c}_5\leq|s_j|\leq u^{(3\tilde{\theta})/4}$, using \eqref{ing104S}, \eqref{zt density 0S}, and the assumption on $m_t(x)$, we have 
\begin{equation*}\label{ptfisherS.3}
\frac{1}{\tilde{\sigma}_j(\mathbf{s})}\phi_j(\frac{s_j}{\tilde{\sigma}_j(\mathbf{s})})\leq \frac{1}{m_j(s_j)}\phi_j(\frac{s_j}{m_j(s_j)})\leq \frac{\bar{c}_3}{\bar{c}_4|s_j|},
\end{equation*}
which implies the desired bound. Finally, for $u^{3\tilde{\theta}/4}<|s_j|$, note that if $\mathbf{s}\in A_{ik,k}(u)$, then $\max\limits_{0\leq j\leq k-1}\tilde{\sigma}^2_{dk^2-j}(\mathbf{s})<u^{\tilde{\theta}}$. 
Combined with \eqref{zt densityS}, this yields
\begin{equation*}\label{ptfisherS.4}
\phi_j(\frac{s_j}{\tilde{\sigma}_j(\mathbf{s})})\leq \phi_j(\frac{s_j}{u^{\frac{\tilde{\theta}}{2}}}),
\end{equation*}
provided $u> (\bar{c}_1\vee\bar{c}_2)^{4/\tilde{\theta}}$.
Additionally, since $\tilde{\sigma}_j(\mathbf{s})\geq c_0$,
the bound
\eqref{ptfisherS.1} follows.
With both 
\eqref{ptfisherS.1} and \eqref{ptfisherS.2} verified,
the proof is complete.
}
\end{proof}

\end{document}
